\subsection{Task 3: Sức chứa theo quận}

\subsubsection{Thiết kế Idiom}

\begin{longtable}{@{}>{\raggedright\arraybackslash}p{0.20\textwidth}>{\raggedright\arraybackslash}p{0.74\textwidth}@{}}
\toprule
\textbf{Đặc điểm} & \textbf{Chi tiết} \\
\midrule
\endhead

\bottomrule
\endfoot

\textbf{Idiom} & Heat Map (Bản đồ nhiệt) \\
\addlinespace[0.08cm]

\textbf{What} & Khu vực: Categorical, Nhóm sức chứa: Ordinal, Số lượng listing: Quantitative \\
\addlinespace[0.08cm]

\textbf{Why} & produce $\rightarrow$ browse, lookup $\rightarrow$ summarize
\begin{itemize}[leftmargin=1.4em, itemsep=0.08em, topsep=0.1cm]
    \item Biểu đồ giúp xác định khu vực nào phù hợp với từng quy mô khách, ví dụ khu vực nào có nhiều listing cho cặp đôi, nhóm nhỏ, gia đình hoặc nhóm đông người.
\end{itemize} \\
\addlinespace[0.08cm]

\textbf{How} & \textbf{Encode:}
\begin{itemize}[leftmargin=1.4em, itemsep=0.08em, topsep=0.1cm]
    \item Mark: Area
    \item Channel:
    \begin{itemize}[leftmargin=1.4em, itemsep=0.08em]
        \item PosX: Nhóm sức chứa
        \item PosY: Khu vực
        \item Hue Color: Số lượng listing
    \end{itemize}
\end{itemize}

\vspace{0.2cm}
\textbf{Manipulate:}
\begin{itemize}[leftmargin=1.4em, itemsep=0.08em, topsep=0.1cm]
    \item Selection: Chọn một ô để highlight tổ hợp khu vực và nhóm sức chứa.
    \item Hover để xem số lượng listing chi tiết.
\end{itemize}

\vspace{0.2cm}
\textbf{Facet:}
\begin{itemize}[leftmargin=1.4em, itemsep=0.08em, topsep=0.1cm]
    \item Có thể facet theo loại phòng.
\end{itemize}

\vspace{0.2cm}
\textbf{Reduce:}
\begin{itemize}[leftmargin=1.4em, itemsep=0.08em, topsep=0.1cm]
    \item Chuẩn hóa sức chứa thành các nhóm: 1--2 khách, 3--4 khách, 5--6 khách, 7+ khách.
\end{itemize} \\
\addlinespace[0.08cm]

\textbf{Scale} &
\begin{itemize}[leftmargin=1.4em, itemsep=0.08em, topsep=0pt]
    \item Main key: khu vực $\times$ nhóm sức chứa
\end{itemize} \\
\end{longtable}

\begin{figure}[H]
    \centering
    \safeincludegraphics[width=0.9\linewidth]{charts/domain_task3_chart2.png}
    \caption{Hình 3.2. Biểu đồ sức chứa theo quận}
    \label{fig:domain-task3-chart2}
\end{figure}

\noindent\small\textit{Xem biểu đồ tương tác: }\href{https://public.tableau.com/app/profile/duy.le4605/viz/Group10_17788537028430/Task3_2-HeatmapCapacitybyBorough?publish=yes}{Tableau Public}\normalsize

\subsubsection{Đánh giá biểu đồ}

\textbf{1. Nguyên lý biểu đạt (Expressiveness):}

\begin{itemize}
    \item Thuộc tính: Khu vực (Categorical -- C), Nhóm sức chứa (Ordinal -- O), Số lượng listing (Quantitative -- Q).
    \item Channel:
    \begin{itemize}
        \item PosX: Nhóm sức chứa (Ordinal).
        \item PosY: Khu vực (Categorical).
        \item Color (Saturation -- Độ bão hòa màu): Thể hiện độ lớn của giá trị định lượng (Số lượng listing).
    \end{itemize}
    \item Đánh giá mức độ quan trọng và phân bổ kênh theo Mackinlay:
    \begin{itemize}
        \item Ưu tiên 1 -- Vị trí không gian (PosX, PosY): Sử dụng hai trục tọa độ để tạo thành lưới ma trận (Matrix) là phương án tối ưu để biểu đạt mối quan hệ giữa hai thuộc tính định danh/thứ tự. Theo Mackinlay, kênh Vị trí là kênh mạnh nhất cho mọi loại dữ liệu, giúp người xem dễ dàng tra cứu (Lookup) tổ hợp giữa Quận và Sức chứa.
        \item Ưu tiên 2 -- Màu sắc (Color Saturation/Luminance): Do hai kênh Vị trí đã bị chiếm dụng cho các chiều phân loại, giá trị định lượng (Q) phải sử dụng kênh màu sắc. Theo thang đo Mackinlay cho dữ liệu định lượng, kênh này xếp hạng thấp hơn Vị trí/Chiều dài nhưng lại cực kỳ hiệu quả trong việc nhận diện các ``điểm nóng'' (Hotspots) hoặc quy luật mật độ trên diện rộng.
    \end{itemize}
    \item Kết luận: Biểu đồ sử dụng chuẩn xác các kênh biểu đạt. Việc phân nhóm (Binning) sức chứa thành 4 cấp độ (1-2, 3-4, 5-6, 7+) giúp tinh gọn dữ liệu và tăng tính biểu đạt cho các User Story về quy mô khách.
\end{itemize}

\textbf{2. Nguyên lý hiệu quả (Effectiveness):}

\begin{itemize}
    \item Độ chính xác (Accuracy):
    \begin{itemize}
        \item Kênh vị trí (PosX, PosY) đảm bảo việc phân loại các nhóm sức chứa và quận huyện đạt độ chính xác tuyệt đối.
        \item Kênh màu sắc (Color Saturation) tuân theo định luật Stevens' Psychophysical law với hàm mũ cảm nhận thấp ($p < 1.0$), dẫn đến sai lệch cảm nhận về lượng (ví dụ: khó phân biệt một ô đậm hơn 10\% có thực sự nhiều hơn 10\% số lượng hay không). Độ lỗi Log\_error ước tính rơi vào khoảng $\approx 2.5$. Tuy nhiên, thiết kế đã khắc phục hoàn toàn nhược điểm này bằng cách hiển thị trực tiếp nhãn số (Labels) bên trong các ô vuông (Mark Area).
    \end{itemize}
    \item Khả năng phân biệt (Discriminability): Dải màu chuyển sắc (Sequential) giúp phân biệt rõ các vùng có mật độ cao (Manhattan 1-2 khách: 5,541) so với các vùng khan hiếm (Staten Island 7+ khách: 16). Khoảng cách giữa các ô (Grid) giúp tách biệt rõ ràng các đơn vị dữ liệu, tránh hiện tượng hòa lẫn màu sắc.
    \item Khả năng tách biệt (Separability): Các kênh vị trí và màu sắc tách biệt tốt. Thao tác Reduce (Chuẩn hóa sức chứa thành các nhóm Ordinal) giúp loại bỏ sự nhiễu loạn của hàng chục giá trị số lẻ, tập trung sự chú ý của người xem vào các phân khúc thị trường chính.
\end{itemize}

\subsubsection{Phân tích biểu đồ}

\begin{itemize}
    \item Phân khúc phòng dành cho 1-2 khách (khách lẻ hoặc cặp đôi) chiếm đa số tuyệt đối tại mọi khu vực, đặc biệt bùng nổ tại Manhattan và Brooklyn.
    \item Nguồn cung tỷ lệ nghịch với sức chứa: số lượng listing giảm mạnh khi quy mô nhóm khách tăng lên, cho thấy Airbnb NYC chủ yếu phục vụ tệp khách nhỏ.
    \item Manhattan và Brooklyn là hai ``điểm nóng'' cung cấp giải pháp lưu trú cho mọi quy mô nhóm khách, trong khi Staten Island gần như đứng ngoài phân khúc nhóm đông người (7+ khách).
    \item Sự tập trung mật độ cao tại ô ``Manhattan / 1-2 khách'' (5,541) xác nhận đây là phân khúc cạnh tranh khốc liệt nhất nhưng cũng có nhu cầu cao nhất toàn thị trường.
\end{itemize}
